% Introduction v1 (20 Nov 2026). Statements about cited papers are abstract-level only (author has not read the full papers).
% Before submission the author must read each cited abstract and confirm. No novelty claim is made (the IGARSS template asks for one; our rules forbid it).

\section{Introduction}
\label{sec:intro}
Subsurface ocean temperature is observed sparsely, whereas satellite sea surface temperature (SST) and sea level are dense in space and time. Estimating the subsurface from the surface is therefore an established problem, and neural networks are increasingly used for it \cite{smith2023}. Recent studies train networks to estimate temperature, and in some cases salinity, from satellite fields: in the tropical Indian Ocean with Argo data \cite{qi2023}, in the Atlantic from in situ vertical profiles \cite{smith2023}, and in the northwestern Pacific from Argo and ship-based profiles \cite{chae2026} or as daily three-dimensional fields from real-time SST and sea level \cite{wang2026}. These studies propose a network and report how accurate it is against observations or other models.

We ask a different question: how much does a deep network add over strong baselines, and how much of the answer depends on the reference it is scored against? We train U-Nets \cite{unet} to reproduce the GLORYS12 reanalysis \cite{glorys12} at 15 depths from 0 to 1000~m over the North Indian Ocean, using only SST \cite{ostia} and sea level anomaly \cite{duacs} as inputs. We compare them with seasonal climatology, a per-cell linear regression and gradient-boosted trees \cite{xgboost} under a protocol fixed before any test period was opened: chronological splits, three test periods, the last of them prospective and pre-registered, and block-bootstrap intervals for every difference. We also score every estimate against Argo profiles. Argo is assimilated by GLORYS12, so the two references are not independent, and we do not present them as such.

This paper makes three contributions. (i)~It measures the gain of a U-Net ensemble over a linear regression against GLORYS12 on three periods, and shows where the gain appears (the thermocline) and where it does not (near the surface and below 300~m). (ii)~It shows that the ranking of models depends on the reference: an ensemble variant that is worst against GLORYS12 is closest to Argo, and a single constant per depth, fitted on validation Argo, recovers essentially the whole difference. (iii)~It states which analyses were pre-registered and which were specified after the test periods were scored. We claim neither real-time use nor skill below 300~m. Section~\ref{sec:data} describes the data and methods, Section~\ref{sec:results} the results and Section~\ref{sec:discussion} their limits.
